% ---
% id: reductions/prop:pinching
% kind: prop
% title: "Sector pinching as a squashing"
% summary: >-
%   Pinching onto finitely many orthogonal sectors is a self-adjoint channel,
%   and a target whose elements are block diagonal for those sectors squashes
%   to itself through it. The pinching can therefore follow a squasher whose
%   target and carried operators are block diagonal.
% prominence: standard
% status: accepted
% uses: [def:squashing]
% mentions: [foundations/def:pinching]
% assumes: []
% refs: []
% provenance: newly-derived
% lean:
% ---

The sector pinching below is the pinching of
\xref{foundations/def:pinching} on the target $\Qsys$: it deletes the
coherences between sectors. The proposition proves the two facts the
squashing constructions use, that it is a self-adjoint channel and that a
block-diagonal target squashes to itself through it.

\begin{proposition}\label{prop:pinching}
Let $\{\Pi_s\}_{s\in\SecSet}$ be finitely many mutually orthogonal projections
on $\Hs{\Qsys}$ with $\sum_s\Pi_s=\idon{\Qsys}$, and let
$\PinS(\sigma)\coloneqq\sum_s\Pi_s\sigma\Pi_s$. Then $\PinS$ is a channel with
$\PinS^{\dagger}=\PinS$, and if every $F_v$ is block diagonal,
$F_v=\sum_s\Pi_sF_v\Pi_s$, then $F\dsq_{\PinS}F$.
\end{proposition}

\begin{proof}
$\PinS$ is CP as a sum of conjugations and trace preserving by cyclicity and
$\sum_s\Pi_s=\idon{\Qsys}$; it is self-adjoint because
$\tr[X\sum_s\Pi_s\sigma\Pi_s]=\tr[(\sum_s\Pi_sX\Pi_s)\sigma]$. Block
diagonality is $\PinS(F_v)=F_v$, which is
$\PinS^{\dagger}(F_v)=F_v$.
\end{proof}
